\begin{pastexam}[label={px:ntust-112-ice-05}]{臺科大資工所 112 學年度 第 5 題}
（複選；三小題）

(15\%) One day, Bob creates three novel sorting algorithms to sort N integer elements stored in an
array ARR. The index begins from 0 for the array, and the pseudo codes are listed below:

\noindent\textbf{Method 1:}
\begin{minted}[fontsize=\small]{text}
method_one (ARR , N)
{
    int TEMP, k, i
    for k=1 to N-1
        TEMP = ARR[k]
        i = k-1
        while TEMP <= ARR[i] && i >= 0
            ARR[i+1] = ARR[i]
            i = i-1
        ARR[i+1] = TEMP
}
\end{minted}

\noindent\textbf{Method 2:}
\begin{minted}[fontsize=\small]{text}
method_two (ARR , N)
{
    int TEMP_a, TEMP_b, k, j
    for k=0 to N-1
        TEMP_a = ARR[k]
        TEMP_b = k
        for j=k+1 to N-1
            if TEMP_a >= ARR[j]
                TEMP_a = ARR[j]
                TEMP_b = j
        TEMP_a = ARR[k]
        ARR[k] = ARR[TEMP_b]
        ARR[TEMP_b] = TEMP_a
}
\end{minted}

\noindent\textbf{Method 3:}
\begin{minted}[fontsize=\small]{text}
method_three (ARR , N)
{
    int TEMP, k, j
    for k=0 to N-1
        for j=0 to N-k
            if ARR[j] >= ARR[j+1]
                TEMP = ARR[j]
                ARR[j] = ARR[j+1]
                ARR[j+1] = TEMP
}
\end{minted}

\begin{enumerate}[label=(\arabic*)]
\item (5\%) Which statement(s) is(are) true about the first method?
\begin{enumerate}[label=(\Alph*)]
\item It does not work for sorting.
\item The best case time complexity is O($n$)
\item The average case time complexity is O($n^2$)
\item It is a stable sorting algorithm
\item None of the above statements is true
\end{enumerate}
\item (5\%) Which statement(s) is(are) true about the second method?
\begin{enumerate}[label=(\Alph*)]
\item The best case time complexity is O($n^2$)
\item The average case time complexity is O($n^2$)
\item The worst case time complexity is O($n^2$)
\item It is a stable sorting algorithm
\item None of the above statements is true
\end{enumerate}
\item (5\%) Which statement(s) is(are) true about the third method?
\begin{enumerate}[label=(\Alph*)]
\item It does not work for sorting.
\item The best case time complexity is O($n^2$)
\item The average case time complexity is O($n$)
\item It is a stable sorting algorithm
\item None of the above statements is true
\end{enumerate}
\end{enumerate}
\end{pastexam}

\begin{pastsoln}
本書的答案：(1) (B)、(C)；(2) (A)、(B)、(C)；(3) (A)、(B)。

三個方法分別是插入排序、選擇排序、氣泡排序（資結書第~14.2~節），但都有改動：比較用了
\texttt{<=} 或 \texttt{>=}，第三個的迴圈範圍也不對。 題目強調「index begins from 0」，下面逐一照
寫出來的程式檢查。 選項裡的 $n$ 就是題目的 $N$。

\textbf{(1) Method 1：插入排序。} 第 $k$ 輪把 \texttt{ARR[k]} 插進已排好的 \texttt{ARR[0..k-1]}。
\begin{itemize}
\item 迴圈條件 \texttt{TEMP <= ARR[i] \&\& i >= 0} 把 \texttt{i >= 0} 寫在後面。 照 C 語言由左到右
的求值，$i = -1$ 時會先讀 \texttt{ARR[-1]}；正確的寫法應該把 \texttt{i >= 0} 放前面。 本書把它當成
虛擬碼的寫法不精確（\texttt{i >= 0} 是迴圈的邊界條件），讀成 $i < 0$ 就停，這樣它能正確排序，
\textbf{(A) 錯}。
\item \textbf{(B) 對。} 輸入已經由小到大排好、而且數字都不同時，每一輪 \texttt{TEMP <= ARR[k-1]}
一開始就不成立，不搬移，$N - 1$ 輪共 $\Theta(N)$。
\item \textbf{(C) 對。} 平均每個數要往前搬過前面約一半的數，總共約 $N^2/4$ 次搬移，$\Theta(N^2)$。
\item \textbf{(D) 錯。} 條件是 \texttt{<=}，碰到相等的數也會繼續往前搬，相等的數順序被顛倒。 例：
$\langle 5_a, 5_b\rangle$，$k = 1$ 時 $\mathit{TEMP} = 5_b \le 5_a$，$5_a$ 往後搬，得
$\langle 5_b, 5_a\rangle$。 標準的插入排序用 \texttt{<}，才是穩定的。
\item (E) 錯，(B)(C) 對。
\end{itemize}

\textbf{(2) Method 2：選擇排序。} 第 $k$ 輪在 \texttt{ARR[k..N-1]} 找最小值的位置
\texttt{TEMP\_b}，再和 \texttt{ARR[k]} 交換。 範圍都正確，能排序。
\begin{itemize}
\item \textbf{(A)(B)(C) 對。} 不管輸入如何，內層迴圈都跑滿：比較次數
$\sum_{k=0}^{N-1} (N - 1 - k) = N(N - 1)/2$，最好、平均、最壞都是 $\Theta(N^2)$，當然也是 $O(N^2)$。
\item \textbf{(D) 錯。} 條件 \texttt{>=} 讓它找到的是\textbf{最後一個}最小值；而且交換會把
\texttt{ARR[k]} 丟到後面去。 例：$\langle 5_a, 5_b\rangle$，$k = 0$ 時 $5_a \ge 5_b$，
$\mathit{TEMP\_b} = 1$，交換得 $\langle 5_b, 5_a\rangle$。 （即使改成 \texttt{>}，選擇排序的長距離交換也不
穩定，例如 $\langle 5_a, 5_b, 3\rangle$。）
\item (E) 錯。
\end{itemize}

\textbf{(3) Method 3：氣泡排序。} 下標從 $0$ 到 $N - 1$，但內層迴圈 $j$ 從 $0$ 跑到 $N - k$，比較
\texttt{ARR[j]} 和 \texttt{ARR[j+1]}：$k = 0$ 時 $j$ 到 $N$，用到 \texttt{ARR[N]}、\texttt{ARR[N+1]}；
$k = 1$ 時 $j$ 到 $N - 1$，用到 \texttt{ARR[N]}。 這兩格在陣列外面。
\begin{itemize}
\item \textbf{(A) 對。} 第一趟把最大值一路換到 \texttt{ARR[N-1]}；接著 $j = N - 1$ 比
\texttt{ARR[N-1] >= ARR[N]}，陣列外那一格的值只要不比最大值大，就會被換進陣列，最大值被換到陣列
外。 結果陣列裡少了原本的數、多了一個不知道是什麼的值，所以照寫出來的程式不能正確排序。 正確的
範圍是 \texttt{j=0 to N-k-2}。
\item \textbf{(B) 對。} 沒有「一趟沒交換就提早結束」的機制，不管輸入如何兩層迴圈都跑滿，最好情況
也是 $\Theta(N^2)$（把範圍改對以後也一樣）。
\item \textbf{(C) 錯。} 平均也是 $\Theta(N^2)$，不是 $O(N)$。
\item \textbf{(D) 錯。} 條件是 \texttt{>=}，相等的相鄰兩數也交換：$\langle 5_a, 5_b\rangle$ 會變成
$\langle 5_b, 5_a\rangle$，不穩定（即使範圍改對也一樣）。
\item (E) 錯。
\end{itemize}

如果把第 (3) 小題的迴圈範圍當成筆誤、只看演算法本身，答案是 (B)。 本書照題目強調的「下標從 $0$
開始」逐字檢查，選 (A)、(B)。
\end{pastsoln}
